% Beweiskarten: sechs je Seite (2 × 3), gestrichelte Schnittlinien. Eine Karte je Beweisstück.
% Die Karten sagen, was die Helden sehen – nicht, welchen Schluss es stützt.
\vspace{-1mm}
{\small\itshape\color{fwgrey}Ausschneiden und den Spielern geben, sobald sie das Beweisstück gefunden haben. Die Nummern
entsprechen der Tabelle im Kapitel \hyperref[beweise]{Die Beweisstücke}.\par}
\vspace{2mm}
\begin{fwcards}
\fwcard{1}{Die Wunden}{Tatort · Fenns Zimmer}{Ein Schlag gegen die Schläfe, eine zweite Wunde am Hinterkopf, Blut am
  Herdstein. Keine Stich- oder Schnittwunde. Der Körper ist kalt und steif: Fenn starb gegen Mitternacht.}
\fwcard{2}{Das Fenster}{Tatort · Fenns Zimmer}{Das Fenster stand offen. Auf der Innenseite der Fensterbank klebt
  Schlamm. Draußen ist das Moos an der Hauswand unberührt, im Beet darunter keine einzige Spur.}
\fwcard{3}{Harnils Stiefel}{Torwächterhütte}{Harnils Stiefel tragen dunkle Walderde, Laub und Kiefernnadeln.
  Kein Krümel roter Lehm.}
\fwcard{4}{Das Wort des Rangers}{Lager im Chetwood · Iorlas}{„Harnil kam kurz vor Mitternacht ins Lager. Er hielt
  Wache bis zur vierten Stunde und brach noch vor Tagesanbruch nach Bree auf.“}
\fwcard{5}{Roter Lehm}{Tatort · Fensterbank}{Der Schlamm auf der Fensterbank ist roter, zäher Lehm. In Bree ist die Erde
  dunkel. „So rot ist’s nur draußen bei Archet und am Rand vom Chetwood.“\par\textit{Mit Explore: Strohhäcksel und Spreu im Lehm – Ackerboden, kein Waldboden.}}
\fwcard{6}{Das Gästebuch}{Schankraum · Butterburs Buch}{Markttag, Obergeschoss:\par
  Zimmer 4 – E. Brambleby, Archet, zum Markt.\par
  Zimmer 5 – O. Fenn, Händler.\par
  Zimmer 6 – zwei Zwerge, beim ersten Licht abgereist.}
\end{fwcards}
\clearpage
\begin{fwcards}
\fwcard{7}{Der hinkende Schatten}{Staddle · Lily Brockhouse}{Kurz nach Mitternacht ging Fenns Fenster auf – niemand
  stieg hinaus oder hinein, Lily und Tom standen darunter. Gegen halb eins kam ein Mann mit Laterne aus Fenns Zimmer,
  ging zur Haupttreppe hin und verschwand in Zimmer 4. Er hinkte.}
\fwcard{8}{Die Stiefel am Trog}{Stallhof · Tom Tunnelly}{Im Morgengrauen wusch ein Mann am Trog im Stallhof roten Lehm
  von seinen Stiefeln: Edric Brambleby aus Archet, der zum Markt im Pony übernachtete.}
\fwcard{9}{Hobs Botgänge}{Bree · Hob Ferny}{Hob trug für Fenn Botschaften zu zwei Fremden, die draußen vor dem Hag
  lagern. Fenn zahlte gut – zu gut für einen Waidhändler. Die Börse hat Hob am Morgen genommen.}
\fwcard{10}{Die Karte aus dem Wagen}{Combe · doppelter Boden}{Eine handgezeichnete Karte des Bree-lands: die Tore, die
  Wege der Ranger über Archet, und am Hag im Nordosten eine Lücke mit den Buchstaben „E. B.“}
\fwcard{11}{Brief an Fenn}{Combe · doppelter Boden}{Ein gefalteter Brief, gezeichnet „R.“, unter der Karte versteckt.
  \par\textit{Dazu: Handout „Brief an Fenn“.}}
\fwcard{12}{Die Lücke im Hag}{Bree · Nordostrand des Hags}{Hinter Brombeeren eine Lücke, gerade groß genug für einen
  Mann. Frische Spuren in südländischen Stiefeln, ältere eines Hinkenden. Diesen Abschnitt bessert Edric Brambleby aus.}
\end{fwcards}
\clearpage
\begin{fwcards}
\fwcard{13}{Lehm in den Nähten}{Archet · Edrics Hof}{Edrics Stiefel stehen frisch geschrubbt neben der Tür. In den
  Nähten sitzt noch roter Lehm.}
\fwcard{14}{Fenns Notizbuch}{Archet · Edrics Werkstatt}{In Öltuch gewickelt, im hohlen Fuß der Hobelbank: das kleine
  Lederbuch, das Fenn am Abend in der Manteltasche trug.\par\textit{Dazu: Handout „Seite aus Fenns Notizbuch“.}}
\fwcard{15}{Das Wort des Südländers}{Ein Gefangener}{„Fenn hatte einen Mann in Archet. Einen, der hinkt. Der hat das
  Buch – und wir wollten es holen.“}
\end{fwcards}
